\documentclass{article}
\usepackage{amsmath,amssymb,amsthm}
\usepackage{algorithm}
\usepackage{algpseudocode}
\usepackage{hyperref}
\newtheorem{theorem}{Theorem}
\newtheorem{lemma}{Lemma}
\newtheorem{assumption}{Assumption}
\newcommand{\prox}{\operatorname{prox}}
\newcommand{\grad}{\nabla}
\newcommand{\inner}[2]{\langle #1,#2\rangle}
\newcommand{\norm}[1]{\|#1\|}
\begin{document}

\section*{Proximal Gradient Method (PGM)}

\section*{Mathematical Formulation}
Let  
\[
F(x)=f(x)+g(x),\qquad x\in\mathbb R^{d},
\]
where  
\begin{itemize}
    \item $f:\mathbb R^{d}\to\mathbb R$ is convex, differentiable and $L$‑smooth:
    \[
    \norm{\grad f(x)-\grad f(y)}\le L\norm{x-y},\quad\forall x,y,
    \]
    \item $g:\mathbb R^{d}\to\mathbb R\cup\{+\infty\}$ is convex (possibly non‑smooth) and its proximal operator
    \[
    \prox_{\eta g}(v):=\arg\min_{u}\Big\{ g(u)+\frac{1}{2\eta}\norm{u-v}^{2}\Big\}
    \]
    is easy to compute.
\end{itemize}
Given a stepsize $\eta\in(0,1/L]$, the proximal gradient iteration is
\begin{equation}
\boxed{%
x_{k+1}= \prox_{\eta g}\bigl(x_{k}-\eta \grad f(x_{k})\bigr),\qquad k=0,1,2,\dots
}
\tag{PGM}
\end{equation}
The algorithm seeks a minimizer $x^{\star}\in\arg\min F$.

\section*{Pseudocode}
\begin{algorithm}[H]
\caption{Proximal Gradient Method}
\begin{algorithmic}[1]
\State \textbf{Input:} stepsize $\eta\in(0,1/L]$, initial point $x_{0}\in\mathbb R^{d}$, max iterations $K$.
\For{$k=0$ \textbf{to} $K-1$}
    \State $v_{k}\gets x_{k}-\eta\,\grad f(x_{k})$ \Comment{gradient step}
    \State $x_{k+1}\gets \prox_{\eta g}(v_{k})$ \Comment{proximal step}
\EndFor
\State \textbf{Output:} $x_{K}$
\end{algorithmic}
\end{algorithm}

\section*{Dry Run Example}
Consider the scalar problem
\[
f(w)=(w-3)^{2},\qquad g\equiv0,
\]
so $\prox_{\eta g}= \operatorname{id}$.  
Take $w_{0}=0$, stepsize $\eta=0.1$, and run three iterations.

\begin{center}
\begin{tabular}{c|c|c|c}
$k$ & $w_{k}$ & $\grad f(w_{k})=2(w_{k}-3)$ & $w_{k+1}=w_{k}-\eta\grad f(w_{k})$ \\ \hline
0 & $0$ & $-6$ & $0-0.1(-6)=0.6$ \\
1 & $0.6$ & $2(0.6-3)=-4.8$ & $0.6-0.1(-4.8)=1.08$ \\
2 & $1.08$ & $2(1.08-3)=-3.84$ & $1.08-0.1(-3.84)=1.464$ \\
\end{tabular}
\end{center}
Objective values:
\[
F(w_{0})=9,\;F(w_{1})=(0.6-3)^{2}=5.76,\;F(w_{2})=(1.08-3)^{2}=3.6864,\;F(w_{3})=(1.464-3)^{2}=2.368.
\]
The objective decreases monotonically, illustrating the method.

\newpage
\section*{Assumptions}
\begin{assumption}[Convexity and Smoothness of $f$]\label{as:fsmooth}
$f$ is convex and $L$‑smooth, i.e.
\[
f(y)\le f(x)+\inner{\grad f(x)}{y-x}+\frac{L}{2}\norm{y-x}^{2},\qquad\forall x,y\in\mathbb R^{d}.
\]
\end{assumption}
\begin{assumption}[Convexity of $g$]\label{as:gconvex}
$g$ is proper, closed and convex.
\end{assumption}
\begin{assumption}[Stepsize]\label{as:stepsize}
The stepsize satisfies $0<\eta\le 1/L$.
\end{assumption}

\section*{Main Convergence Theorem}
\begin{theorem}[Ergodic $O(1/k)$ rate]\label{thm:ergodic}
Let $\{x_{k}\}$ be generated by \eqref{PGM} under Assumptions~\ref{as:fsmooth}–\ref{as:stepsize}, and let $x^{\star}\in\arg\min F$. Then for any $K\ge1$,
\[
\frac{1}{K}\sum_{k=0}^{K-1}\bigl(F(x_{k})-F(x^{\star})\bigr)
\;\le\;
\frac{\norm{x_{0}-x^{\star}}^{2}}{2\eta K}.
\tag{1}
\]
Consequently,
\[
\min_{0\le k\le K-1}\bigl(F(x_{k})-F(x^{\star})\bigr)\le
\frac{\norm{x_{0}-x^{\star}}^{2}}{2\eta K}.
\]
\end{theorem}

\section*{Theoretical Properties}
\begin{itemize}
    \item \textbf{Stability:} The iteration is a contraction of the composite map
    $T_{\eta}:= \prox_{\eta g}\circ(\operatorname{id}-\eta\grad f)$ because both components are non‑expansive (the proximal operator) or firmly non‑expansive (gradient step with $\eta\le1/L$). Hence the sequence $\{x_{k}\}$ remains bounded.
    \item \textbf{Hyperparameter Sensitivity:} The bound \eqref{1} improves linearly with $\eta$; the maximal admissible stepsize $1/L$ yields the smallest constant $\frac{1}{2\eta}$.
    \item \textbf{Comparison:} Compared with plain gradient descent (where $g\equiv0$) the rate is identical, but PGM additionally handles the non‑smooth term $g$ without sacrificing the $O(1/k)$ ergodic rate.
\end{itemize}

\section*{Discussion}
The theorem guarantees that the average objective gap decays as $1/K$. In practice the last iterate often enjoys the same rate under stronger assumptions (e.g. strong convexity). The proof below is fully self‑contained and relies only on the non‑expansiveness of the proximal operator and the smoothness inequality for $f$.

\newpage
\section*{Preliminaries (Definitions and Lemmas)}
\begin{lemma}[Non‑expansiveness of the proximal operator]\label{lem:prox_ne}
For any $\eta>0$ and any $u,v\in\mathbb R^{d}$,
\[
\norm{\prox_{\eta g}(u)-\prox_{\eta g}(v)}\le\norm{u-v}.
\tag{2}
\]
\end{lemma}
\begin{proof}
The proximal operator is the resolvent of the monotone sub‑differential $\partial g$, and resolvents are firmly non‑expansive. Hence (2) holds. \qedhere
\end{proof}

\begin{lemma}[Descent lemma for $L$‑smooth convex $f$]\label{lem:descent}
For all $x,y\in\mathbb R^{d}$,
\[
f(y)\le f(x)+\inner{\grad f(x)}{y-x}+\frac{L}{2}\norm{y-x}^{2}.
\tag{3}
\]
\end{lemma}
\begin{proof}
Standard consequence of $L$‑Lipschitz gradient; see Assumption~\ref{as:fsmooth}. \qedhere
\end{proof}

\section*{Main Proof (Step‑by‑Step Derivation)}
We prove Theorem~\ref{thm:ergodic}. Throughout, fix an arbitrary optimal point $x^{\star}\in\arg\min F$.

\noindent\textbf{[Step 1]}  Optimality condition of the proximal step.  
From the definition of $\prox_{\eta g}$, $x_{k+1}$ satisfies
\[
\frac{1}{\eta}\bigl(x_{k}-\eta\grad f(x_{k})-x_{k+1}\bigr)\in\partial g(x_{k+1}).
\tag{4}
\]

\noindent\textbf{[Step 2]}  Subgradient inequality for $g$.  
Since $g$ is convex,
\[
g(x_{k+1})-g(x^{\star})\le
\inner{s_{k+1}}{x_{k+1}-x^{\star}},\qquad
s_{k+1}:=\frac{1}{\eta}\bigl(x_{k}-\eta\grad f(x_{k})-x_{k+1}\bigr).
\tag{5}
\]

\noindent\textbf{[Step 3]}  Combine (5) with the definition of $s_{k+1}$:
\[
g(x_{k+1})-g(x^{\star})\le
\frac{1}{\eta}\inner{x_{k}-\eta\grad f(x_{k})-x_{k+1}}{x_{k+1}-x^{\star}}.
\tag{6}
\]

\noindent\textbf{[Step 4]}  Expand the inner product in (6):
\[
\inner{x_{k}-x_{k+1}}{x_{k+1}-x^{\star}}
-\eta\inner{\grad f(x_{k})}{x_{k+1}-x^{\star}}.
\tag{7}
\]

\noindent\textbf{[Step 5]}  Use the identity
\[
2\inner{a-b}{b-c}= \norm{a-c}^{2}-\norm{a-b}^{2}-\norm{b-c}^{2},
\]
with $a=x_{k}$, $b=x_{k+1}$, $c=x^{\star}$, to rewrite the first term of (7):
\[
2\inner{x_{k}-x_{k+1}}{x_{k+1}-x^{\star}}
= \norm{x_{k}-x^{\star}}^{2}-\norm{x_{k}-x_{k+1}}^{2}-\norm{x_{k+1}-x^{\star}}^{2}.
\tag{8}
\]

\noindent\textbf{[Step 6]}  Divide (8) by $2\eta$ and plug into (6):
\[
g(x_{k+1})-g(x^{\star})\le
\frac{1}{2\eta}\bigl(\norm{x_{k}-x^{\star}}^{2}
-\norm{x_{k}-x_{k+1}}^{2}
-\norm{x_{k+1}-x^{\star}}^{2}\bigr)
-\inner{\grad f(x_{k})}{x_{k+1}-x^{\star}}.
\tag{9}
\]

\noindent\textbf{[Step 7]}  Add $f(x_{k+1})-f(x^{\star})$ to both sides of (9).  
By convexity of $f$,
\[
f(x_{k+1})-f(x^{\star})\le
\inner{\grad f(x_{k})}{x_{k+1}-x^{\star}}
+\frac{L}{2}\norm{x_{k+1}-x_{k}}^{2},\qquad\text{(Lemma~\ref{lem:descent})}
\tag{10}
\]
where we have used $y=x_{k+1}$ and $x=x_{k}$.

\noindent\textbf{[Step 8]}  Adding (9) and (10) cancels the linear term
$\inner{\grad f(x_{k})}{x_{k+1}-x^{\star}}$:
\[
\begin{aligned}
F(x_{k+1})-F(x^{\star})
&\le
\frac{1}{2\eta}\bigl(\norm{x_{k}-x^{\star}}^{2}
-\norm{x_{k+1}-x^{\star}}^{2}\bigr)
+\Bigl(\frac{L}{2}-\frac{1}{2\eta}\Bigr)\norm{x_{k+1}-x_{k}}^{2}.
\end{aligned}
\tag{11}
\]

\noindent\textbf{[Step 9]}  By Assumption~\ref{as:stepsize}, $\eta\le 1/L$, hence
\[
\frac{L}{2}-\frac{1}{2\eta}\le0.
\]
Therefore the second term in (11) is non‑positive and can be dropped:
\[
F(x_{k+1})-F(x^{\star})
\le
\frac{1}{2\eta}\Bigl(\norm{x_{k}-x^{\star}}^{2}
-\norm{x_{k+1}-x^{\star}}^{2}\Bigr).
\tag{12}
\]

\noindent\textbf{[Step 10]}  Rearranging (12) yields a one‑step descent inequality:
\[
\norm{x_{k+1}-x^{\star}}^{2}
\le
\norm{x_{k}-x^{\star}}^{2}
-2\eta\bigl(F(x_{k+1})-F(x^{\star})\bigr).
\tag{13}
\]

\noindent\textbf{[Step 11]}  Sum (13) for $k=0,\dots,K-1$:
\[
\norm{x_{K}-x^{\star}}^{2}
\le
\norm{x_{0}-x^{\star}}^{2}
-2\eta\sum_{k=0}^{K-1}\bigl(F(x_{k+1})-F(x^{\star})\bigr).
\tag{14}
\]

\noindent\textbf{[Step 12]}  Dropping the non‑negative term $\norm{x_{K}-x^{\star}}^{2}$ and shifting the index $k\leftarrow k+1$ give
\[
\sum_{k=0}^{K-1}\bigl(F(x_{k})-F(x^{\star})\bigr)
\le
\frac{\norm{x_{0}-x^{\star}}^{2}}{2\eta}.
\tag{15}
\]

\noindent\textbf{[Step 13]}  Divide (15) by $K$ to obtain the ergodic bound \eqref{1}:
\[
\frac{1}{K}\sum_{k=0}^{K-1}\bigl(F(x_{k})-F(x^{\star})\bigr)
\le
\frac{\norm{x_{0}-x^{\star}}^{2}}{2\eta K}.
\tag{16}
\]

\noindent\textbf{[Step 14]}  Since the minimum of a set is bounded by its average,
\[
\min_{0\le k\le K-1}\bigl(F(x_{k})-F(x^{\star})\bigr)
\le
\frac{1}{K}\sum_{k=0}^{K-1}\bigl(F(x_{k})-F(x^{\star})\bigr),
\]
the second inequality of Theorem~\ref{thm:ergodic} follows directly from (16). \qed

\section*{Rate Derivation (With Constants)}
From \eqref{16} we can read the explicit constant:
\[
\boxed{%
F(x_{k})-F(x^{\star})\le \frac{\norm{x_{0}-x^{\star}}^{2}}{2\eta k},
\qquad k\ge1.}
\]
If we set the maximal stepsize $\eta=1/L$, the bound becomes
\[
F(x_{k})-F(x^{\star})\le \frac{L\,\norm{x_{0}-x^{\star}}^{2}}{2k}.
\]

\section*{Special Cases}
\begin{itemize}
    \item \textbf{Pure gradient descent ($g\equiv0$).}  $\prox_{\eta g}=\operatorname{id}$ and the proof reduces to the classic $O(1/k)$ rate for smooth convex optimization.
    \item \textbf{Indicator function $g=\iota_{\mathcal C}$ of a closed convex set $\mathcal C$.}  The proximal step becomes Euclidean projection onto $\mathcal C$, and the same rate holds for the projected gradient method.
    \item \textbf{Strongly convex $F$.}  If $F$ is $\mu$‑strongly convex, a simple modification of Step 9 yields a linear rate
    $F(x_{k})-F(x^{\star})\le \bigl(1-\eta\mu\bigr)^{k}\bigl(F(x_{0})-F(x^{\star})\bigr)$.
\end{itemize}

\end{document}